\documentclass[11pt]{article}
\usepackage{Style}

\nocite{*} % Comentar si quiero citar
%\addbibresource{bibliografia.bib} % Quitar el comentado si quiero usar bibliografia

\begin{document}

\begin{titlepage} % Portada con información necesaria
    \hspace*{-2cm}
    \raisebox{-2cm}[0pt][0pt]{
        \includegraphics[width=0.4\textwidth]{logo_universidad2.png}
    }

    \vspace{3.5cm}
    \begin{center}
        {\Large Pontificia Universidad Católica de Chile}
        
        \vspace{2cm}
        {\Huge \bfseries 
        Extensiones de Grupos
        }
        
        \vspace{1.5cm}
        {\Large Benjamín Maldonado, Benjamín Mateluna, Sergio Peña}
        
        \vspace{1.5cm}
        {\Large Docente: Sebastián Barbieri}
        \vfill
        
        {\large \today \par}
    \end{center}
\end{titlepage}

\tableofcontents % Índice de contenidos

\vspace{5cm}
\section*{Introducción}
\addcontentsline{toc}{section}{Introducción}

\newpage
\section{Secuencias Exactas Cortas}
\noindent En esta sección estudiaremos la noción de secuencia exacta corta, además, veremos que 
el producto directo y semidirecto se dan cuando la secuencia escinde por izquierda y por derecha 
respectivamente. La idea es codificar la información de dichas estructuras en un objeto más 
general como las secuencias exactas cortas. Comenzamos dando la definición formal de 
esta última %(Modificar)

\begin{dfn}
    Sean $\varphi:A\to B$ y $\psi:B\to C$ morfismos de grupos, decimos que forman una secuencia 
    exacta corta si $\varphi$ es inyectivo, $\psi$ sobreyectivo y $\im{\varphi}=\kr{\psi}$ y lo 
    anotamos

    \vspace{2mm}
    \centerline{
        \xymatrix{
            0 \ar[r] & A \ar[r]^{\varphi} & B \ar[r]^{\psi} & C \ar[r] & 0
        }
    }
\end{dfn}

\vspace{2mm}

\noindent Este tipo de estruturas son naturales y surgen de estudiar el primer teorema de 
isomorfismo. En efecto, como $\kr{\psi}=\im{\varphi}$, tenemos que $\varphi(A)\unlhd B$ y como 
$\psi$ es sobreyectiva concluimos que
\begin{equation*}
    B/\varphi(A)\simeq C
\end{equation*}
De este modo, dado que $\varphi$ es inyectivo pensamos en $A$ como un subgrupo normal de $B$ y en 
$C$ como el grupo cociente $B/A$.

\begin{ejm}
    Dado $n\in\N$, la siguiente secuencia

    \vspace{2mm}
    \centerline{
        \xymatrix{
            0 \ar[r] & n\Z \ar[r]^{i} & \Z \ar[r]^-{\pi} & \Z/n\Z \ar[r] & 0
        }
    } \vspace{2mm}

    es una secuencia exacta corta. 
\end{ejm}

\obs Notemos que si $B$ es un grupo simple entonces la secuencia es trivial en el siguiente 
sentido; como $B$ es simple y $\varphi(A)$ un subgrupo normal isomorfo a $A$, resulta que 
$A\simeq B$ ó $A\simeq1$ y en este último, por exactitud, se tiene que $C\simeq B$.

\vspace{2mm}
\noindent Sabemos que $B/A\simeq C$, pero la pregunta interesante es ¿Cuándo podemos obtener 
información sobre la estructura de $B$ a partir de $A$ y $C$? ¿Que propiedades debe cumplir la 
secuencia exacta? Supongamos que $A\times C\simeq B$, consideremos los morfismos

\centerline{
    \xymatrix @R=0.5mm @C=2mm{
        i:A\to A\times C & \hhtext{y} & \pi:A\times C\to C \\
        \hspace{3mm} a\mapsto(a,1) & & \hspace{6mm} (a,c)\mapsto c
    }
} \vspace{2mm}

\noindent Notemos que $\im{i}=\kr{\pi}$. Si $\varphi:A\times C\to B$ es isomorfismo, entonces

\centerline{
        \xymatrix @C=1cm{
            0 \ar[r] & A \ar[r]^-{\varphi\circ i} 
            & B \ar[r]^-{\pi\circ\varphi^{-1}} & C \ar[r] & 0
        }
    } \vspace{2mm}

\noindent es una secuencia exacta corta, además, la función $j:A\times C\to A$ dada por $j(a,c)=a$ 
es un morfismo de grupos tal que $j\circ i=id_{A}$. Lo anterior da pie a la siguiente definición

\begin{dfn}
    Decimos que una secuencia exacta corta

    \vspace{2mm}
    \centerline{
        \xymatrix{
            0 \ar[r] & A \ar[r]^{\varphi} & B \ar[r]^{\psi} & C \ar[r] & 0
        }
    } \vspace{2mm}

    se escinde por izquierda, si existe un morfismo de grupos $\phi:B\to A$ tal que 
    $\phi\circ\varphi=id_{A}$.
\end{dfn}

\noindent Con la discusión anterior probamos que si $A\times C\simeq B$ entonces existe una 
secuencia exacta corta que se escinde por izquierda ¿El reciproco será cierto? Por fortuna, el
reciproco también es cierto,

\begin{lema}
    Sean $A, B$ y $C$ grupos. Entonces existe una secuencia exacta corta

    \centerline{
        \xymatrix{
            0 \ar[r] & A \ar[r]^{\varphi} & B \ar[r]^{\psi} & C \ar[r] & 0
        }
    } \vspace{2mm}
    que se escinde por izquierda si y solo si $A\times C\simeq B$.
\end{lema}

\newpage
\begin{proof}
    Ya probamos que $A\times C\simeq B$ entonces existe una secuencia exacta corta que se escinde 
    por izquierda. Por otro lado, supongamos que tal secuencia existe, 
    \begin{equation*}
        \phi:B\to A \hhtext{morfismo tal que } \phi\circ\varphi=id_{A}
    \end{equation*}
    Puesto que la secuencia es exacta se sigue que $\varphi(A)\unlhd B$. Afirmamos que 
    $\varphi(A)\cap\kr{\phi}=\{1\}$ y $\varphi(A)\kr{\phi}$ y por lo tanto 
    $\varphi(A)\times\kr{\phi}\simeq B$. Sea $b\in B$, notemos que
    \begin{equation*}
        b=\varphi\phi(b)(\varphi\phi(b)^{-1}b)
    \end{equation*}
    Claramente el primer término pertenece a $\varphi(A)$, mientras que 
    $\varphi\phi(b)^{-1}b\in\kr{\phi}$, pues $\phi(\varphi\phi(b)^{-1}b)=\phi(b)^{-1}\phi(b)=1$. 
    Supongamos que $b\in \varphi(A)\cap\kr{\phi}$, entonces
    \begin{equation*}
        b=\varphi(a) \htext{para algún }a\in A
    \end{equation*}
    aplicando $\phi$ observamos que $1=\phi(b)=\phi\varphi(a)=a$, es decir, $b=1$. Hemos probado
    que $\varphi(A)\times\kr{\phi}\simeq B$, para finalizar notamos que $\varphi(A)\simeq A$, ya 
    que, $\varphi$ es inyectivo; resta ver que $\kr{\phi}\simeq C$. Veamos que 
    $\psi:\kr{\phi}\to C$ es isomorfismo. Existe $b\in B$ tal que $\psi(b)=c$, entonces
    \begin{equation*}
        c=\psi(b)=\psi(\varphi\phi(b)^{-1}b)
    \end{equation*}
    lo que prueba la sobreyectividad. Para la inyectividad observamos lo siguiente
    $\kr{\phi}\cap\kr{\psi}=\kr{\phi}\cap\im{\varphi}=\{1\}$. Concluimos que $A\times C\simeq B$
    como queriamos.
\end{proof}

\noindent Hemos logrado caracterizar cuando un grupo es isomorfo al producto directo de dos 
grupos; lo que nos permite, por ejemplo, descartar cuando un grupo no corresponde al producto
directo de dos grupos, veremos un ejemplo de esto más adelante. Nuestro siguiente objetivo es el
producto semidirecto, para ello definimos lo siguiente

\begin{dfn}
    Decimos que una secuencia exacta corta
    
    \centerline{
        \xymatrix{
            0 \ar[r] & A \ar[r]^{\varphi} & B \ar[r]^{\psi} & C \ar[r] & 0
        }
    } \vspace{2mm}
    se escinde por derecha, si existe $\phi:C\to B$ morfismo de grupos tal que 
    $\psi\circ\phi=id_{C}$.
\end{dfn}

\noindent Y el siguiente lema nos garantiza un resultado similar en el caso del producto 
semidirecto,

\begin{lema}[Splitting lemma]
    Sean $A, B$ y $C$ grupos. Entonces existe una secuencia exacta corta

    \centerline{
        \xymatrix{
            0 \ar[r] & A \ar[r]^{\varphi} & B \ar[r]^{\psi} & C \ar[r] & 0
        }
    } \vspace{2mm}
    que se escinde por derecha si y solo si $A\rtimes_{\alpha} C\simeq B$.
\end{lema}

\begin{proof}
    Supongamos que $A\rtimes_{\alpha} C\simeq B$, queremos probar que existe una secuencia exacta 
    corta que se escinde por derecha. Al igual que en el caso del producto directo, consideramos los 
    morfismos
    
    \centerline{
        \xymatrix @R=0.5mm @C=2mm{
            i:A\to A\rtimes_{\alpha} C & \hhtext{y} & \pi:A\rtimes_{\alpha} C\to C \\
            \hspace{3mm} a\mapsto(a,1) & & \hspace{6mm} (a,c)\mapsto c
        }
    } \vspace{2mm}
    Aunque en este hay que tener cuidado con la operación en $A\rtimes_{\alpha} C$. Además, 
    consideramos la función $\rho:C\to A\rtimes_{\alpha} C$ dada por $\rho(c)=(1,c)$, y para ver
    es morfismo observamos lo siguiente 
    $\rho(c)\rho(c')=(1,c)(1,c')=(1\alpha_{c}(1),cc')=\rho(cc')$.
    
    \vspace{1mm}
    Es facil ver que $\pi\circ\rho=id_{C}$. Luego, sea $\varphi:A\rtimes_{\alpha}C\to B$ 
    isomorfismo, entonces

    \centerline{
        \xymatrix @C=1cm{
            0 \ar[r] & A \ar[r]^{\varphi\circ i} & B \ar[r]^{\pi\circ\varphi^{-1}} & C \ar[r] & 0
        }
    } \vspace{2mm}
    es una secuencia exacta corta que se escinde por derecha.

    \newpage
    \vspace{2mm}
    Supongamos que la secuencia exacta corta se escinde por derecha, es decir, existe 
    $\phi:C\to B$ morfismo de grupos tal que $\psi\circ\phi=id_{C}$. Como $\varphi$ es un morfismo 
    inyectivo, $A\simeq \varphi(A)$, más aún, dado que la secuencia es exacta 
    $\varphi(A)=\kr{\psi}\unlhd B$ y dado que $\phi$ posee inversa por izquierda se tiene que
    $C\simeq \phi(C)$.

    \vspace{2mm}
    Afirmamos que $\varphi(A)\cap\phi(C)=\{1\}$ y $\varphi(A)\phi(C)=B$, en dicho caso, por el 
    teorema del producto semidirecto concluimos. En efecto, sea $b\in \varphi(A)\cap\phi(C)$, 
    entonces
    \begin{equation*}
        b=\varphi(a) \htext{para algún }a\in A
    \end{equation*}
    así, $\psi(b)=\psi(\varphi(a))=1$. Por otro lado, existe $c\in C$ tal que $b=\phi(c)$, usando
    lo anterior vemos que
    \begin{equation*}
        c=\psi(\phi(c))=\psi(b)=1
    \end{equation*}
    Concluimos que $b=1$. Sea $b\in B$, digamos $c=\psi(b)$ y notemos que $\psi(b(\phi(c^{-1})))
    =cc^{-1}=1$, luego, por exactitud existe $a\in A$ tal que $\varphi(a)=b\phi(c^{-1})$. Por lo 
    tanto,
    \begin{equation*}
        b=\varphi(a)\phi(c)
    \end{equation*}
\end{proof}

\subsection{Ejemplos}

\noindent Notemos que el producto directo en particular es un producto semidirecto inducido por el
morfismo trivial de $C$ en $aut(A)$ y por lo tanto toda secuencia exacta que se escinde por 
izquierda también se escinde por derecha, sin embargo, el recíproco no es cierto, para ello veamos
el siguiente ejemplo.

\begin{ejm}
    Sea $D_{6}$ el grupo dihedral de $6$ elementos (el grupo que corresponde a las simetrías del
    triángulo). Consideremos el sugbrupo normal $A=\{id, r, r^{2}\}$, es normal puesto que su 
    índice es dos. Además,
    \begin{equation*}
        D_{6}/A\simeq\Z_{2}\simeq C:=\{id,s\}
    \end{equation*}
    Luego, la siguiente secuencia exacta
    
    \vspace{2mm}
    \centerline{
        \xymatrix{
            0 \ar[r] & A \ar[r]^{i} & D_{6} \ar[r]^{\pi} & C \ar[r] & 0
        }
    } \vspace{2mm}
    \noindent no escinde por izquierda, de lo contrario $D_{6}\simeq\Z_{6}$, pero sabemos que el
    grupo dihedral de $6$ elementos no es abeliano. Por otro lado, por el teorema de producto 
    semidirecto, resulta que
    \begin{equation*}
        D_{6}\simeq A\rtimes C
    \end{equation*}
    lo que implica que la secuencia escinde por derecha.
\end{ejm}

\noindent Otro ejemplo interesante son los cuaterniones un grupo no abeliano, dado por la 
siguiente presentación
\begin{equation*}
    Q_{8}:=\gen{-1,i,j,k|(-1)^{2}=1,i^{2}=j^{2}=k^{2}=ijk=-1}
\end{equation*}

\begin{ejm}
    No existe una secuencia exacta
    
    \vspace{2mm}
    \centerline{
        \xymatrix{
            0 \ar[r] & A \ar[r]^{\varphi} & Q_{8} \ar[r]^{\psi} & C \ar[r] & 0
        }
    } \vspace{2mm}
    \noindent que se escinda por izquierda o por derecha. En otras palabras, los cuaterniones no
    corresponden ni a un producto directo ni semidirecto de grupos.
    (Pendiente)
\end{ejm}

\newpage
\section{Grupos Solubles y Residualmente finitos}

\begin{prop}
    Sea

    \centerline{
        \xymatrix{
            0 \ar[r] & A \ar[r]^{\varphi} & B \ar[r]^{\psi} & C \ar[r] & 0
        }
    } \vspace{2mm}
    una secuencia exacta corta tal que $A$ es soluble y $C$ es soluble, entonces $B$ es soluble.
\end{prop}

%\printbibliography % Quitar el comentado si quiero usar bibliografia

\end{document}